\documentclass[10pt,a4paper]{amsart}
\usepackage[margin=19mm]{geometry}
\usepackage{amsmath,amssymb,mathtools}
\usepackage[scr=rsfs,cal=euler]{mathalfa}
\usepackage{xfrac}
\usepackage[unicode,hidelinks,bookmarks=false]{hyperref}
\newcommand{\ie}{i.e.\ }
\newcommand{\cf}{cf.\ }
\newcommand{\resp}{respectively\ }
\newcommand{\Subsection}{\subsection}
\providecommand{\nobreakdash}{\nobreak\hbox{-}}
\setlength{\emergencystretch}{2em}
\multlinegap0pt
\usepackage{amssymb}
\newtheorem{theorem}{Theorem}
\newtheorem{corollary}{Corollary}
\newtheorem{proposition}{Proposition}
\def\Z{\mathbb{Z}}
\title{Differences of two squares of upper-triangular $2\times 2$ integer matrices}
\author{Andrej Dujella, Zrinka Franu\v{s}i\'c}
\date{Source: arXiv:2604.23404v2 (CC BY 4.0)}
\begin{document}
\maketitle

\noindent\emph{Source and licence.} Differences of two squares of upper-triangular $2\times 2$ integer matrices. Version arXiv:2604.23404v2 (CC BY 4.0), distributed by arXiv under the Creative Commons Attribution 4.0 International licence, http://creativecommons.org/licenses/by/4.0/, as recorded in the arXiv licence metadata for this exact version and shown on the abstract page https://arxiv.org/abs/2604.23404v2. Full licence notice: \texttt{licenses/arxiv-2604.23404-CC-BY-4.0.txt}.\par\smallskip

\noindent\textit{Editorial scope.} Counted unit: the article's proposition prop:4k-4n-complete (source0.tex lines 372--381). The article's complete original proof of it is delivered with it (source0.tex lines 382--439, one \texttt{proof} environment of 58 lines). No mathematics is written, repaired or re-derived: the statement and the proof are the article's own lines, in the article's own notation, and no retained line is substituted or re-flowed.  Retained uncounted as the source-local context the proof needs: lines 105--114; lines 357--368.  Nothing was omitted from any retained span, so the package carries no omission record.  No retained line contains a citation, so the wrapper carries no bibliography and no bibliography entry.  No retained line contains a figure environment, an image-inclusion command or drawing code, so no image is delivered and the package declares no asset dependency.  Editorial formatting only, in addition: an AMS \texttt{amsart} preamble that restates the article's own declarations and macros used by the retained lines, namely the environments \texttt{theorem}, \texttt{corollary}, \texttt{proposition} on separate counters, together with the standard packages those lines need.  The retained environments print in this document as Theorem 1, Corollary 1, Proposition 1 in order of appearance, whereas the article prints the same items with the numbering of its own full document.  The complete native source of the article as distributed in the arXiv e-print for this exact version is bundled unchanged as \texttt{sources/199163/source0.tex}.
\begin{theorem}\label{thm:main-criterion}
The upper-triangular matrix
$T=\begin{pmatrix}p&r\\0&q\end{pmatrix}$ in $UT_2(\Z)$
can be represented as a difference of two squares in $UT_2(\Z)$  if and only if there exist integers $a,x,d,u$ such that
$$p=a^2-x^2,\qquad q=d^2-u^2,$$
and either
$$(a+d,x+u)\neq(0,0)\quad\text{and}\quad \gcd(a+d,x+u)\mid r,$$
or
$$a+d=0,\quad x+u=0,\quad r=0.$$
\end{theorem}

\begin{corollary}\label{cor:nonrep-mod16}
Let
\[
M=\begin{pmatrix}4i&r\\0&4j\end{pmatrix}\in UT_2(\mathbb Z_{16}),
\qquad i,j\in\{0,1,2,3\}.
\]
Then $M$ cannot be represented as a difference of two squares in $UT_2(\mathbb Z_{16})$ if and only if one of the following holds:
\begin{enumerate}
\item$(i,j)\in\{(1,1),(3,3)\}$ and $r\not\equiv 0\pmod 4$;
\item$(i,j)\in\{(1,3),(2,2),(3,1)\}$ and $r\not\equiv 0\pmod 2$.
\end{enumerate}
\end{corollary}

\begin{proposition}\label{prop:4k-4n-complete}
Let
$$M=\begin{pmatrix}4k&r\\0&4n\end{pmatrix}\in UT_2(\Z).$$
Then $M$ is representable as a difference of two squares in $UT_2(\Z)$ if and only if one of the following conditions holds:
\begin{enumerate}
\item $k+n\equiv1\pmod 2$, or $(k,n)\bmod 4\in\{(0,0),(0,2),(2,0)\}$, with arbitrary $r\in\mathbb Z$;
\item $k+n\equiv0\pmod 4$, $(k,n)\bmod 4\neq(0,0)$, and $r\equiv0\pmod 2$;
\item $k+n\equiv2\pmod 4$,$(k,n)\bmod 4\notin\{(0,2),(2,0)\}$, and $r\equiv0\pmod 4$.
\end{enumerate}
\end{proposition}

\begin{proof}  $\boxed{\Leftarrow}$:
 First, the choice
$$a=k+1,\qquad x=k-1,\qquad d=n+1,\qquad u=n-1,$$
shows representability whenever $r$ satisfies the divisibility condition determined by $\gcd(k+n+2,4)$.
Indeed,
$$a^2-x^2=4k,\qquad d^2-u^2=4n,$$
and
$$a+d=k+n+2,\qquad x+u=k+n-2.$$
Hence
$$\gcd(a+d,x+u)=\gcd(k+n+2,k+n-2)=\gcd(k+n+2,4).$$
Therefore:
\begin{enumerate}
\item[(a)] if $k+n\equiv 1 \pmod 2$, then $\gcd(a+d,x+u)=1$;
\item[(b)] if $k+n\equiv 0 \pmod 4$, then $\gcd(a+d,x+u)=2$;
\item[(c)] if $k+n\equiv 2 \pmod 4$, then $\gcd(a+d,x+u)=4$.
\end{enumerate}
In each case, the assumed divisibility condition on $r$ implies that
$\gcd(a+d,x+u)\mid r.$
Hence Theorem~\ref{thm:main-criterion} shows that $M$ is representable as a difference of two squares in $UT_2(\Z)$ in the following cases:
\begin{itemize}
   \item $k+n\equiv1\pmod 2$,
\item $k+n\equiv0\pmod 4$ and $r\equiv0\pmod 2$;
\item $k+n\equiv2\pmod 4$ and $r\equiv0\pmod 4$.
\end{itemize}

For the remaining cases, we also give concrete representations.  For the case $k\equiv n\equiv 0 \pmod 4$, we put $k=4K$ and $n=4N$. For
$$A=\begin{pmatrix}2(K+1)&b\\0&-(4N+1)\end{pmatrix},\qquad
B=\begin{pmatrix}2(K-1)&y\\0&-(4N-1)\end{pmatrix},$$
we obtain
$$A^2-B^2=\begin{pmatrix}
16K & b(2K-4N+1)-y(2K-4N-1)\\0 & 16N
\end{pmatrix}.$$
Since
$$\gcd(2K-4N+1,\;2K-4N-1)=1,$$
the upper-right entry can be chosen arbitrarily. Hence every matrix
$$\begin{pmatrix}16K&r\\0&16N\end{pmatrix}$$
is representable as a difference of two squares in $UT_2(\Z)$.


Assume $k\equiv 0 \pmod 4$ and $n\equiv 2 \pmod 4$, say $k=4K$ and $n=4N+2$. Consider
$$A=\begin{pmatrix}2(K+1)&b\\0&-(4N+3)\end{pmatrix},\qquad
B=\begin{pmatrix}2(K-1)&y\\0&4N+1\end{pmatrix}.$$
Then
$$A^2-B^2=\begin{pmatrix}
16K &
b(2K-4N-1)+y(2K-4N-3)\\0 &8(2N+1)\end{pmatrix}.$$
Since
$$\gcd(2K-4N-1,\;2K-4N-3)=1,$$
the upper-right entry can be chosen arbitrarily. Hence every matrix of the form
$$\begin{pmatrix}
16K&r\\0&8(2N+1)
\end{pmatrix}$$
is representable as a difference of two squares in $UT_2(\mathbb Z)$. The case $(k,n)\equiv(2,0)\pmod 4$ is analogous.


$\boxed{\Rightarrow}$:
The claim follows from Corollary~\ref{cor:nonrep-mod16}, by contraposition of the fact that representability over $\mathbb Z$ implies representability modulo $16$.
\end{proof}

\end{document}
